\documentclass{scrartcl}
\usepackage[utf8]{inputenc}
\usepackage{fancyhdr}
\usepackage{graphicx}
\usepackage{dirtytalk}
\usepackage[portrait, margin=1in]{geometry}
\usepackage{amsmath, amssymb}
\setlength{\parskip}{1em}
\setlength{\parindent}{0cm} 
\usepackage[mdthm,simplethm]{test}
\usepackage[utf8]{inputenc}
\usepackage{float}
\usepackage{physics}
\setlength{\parindent}{0pt}

\title{Lecture 1}
\author{albert ye}
\date{May 2022}

\begin{document}
\maketitle
\section{Grading}
\begin{enumerate}
	\item Homework: 0\%
	\item Quizzes: 20\%
	\item Participation: 10\%
	\item Midterms: 15\% each, 2 Midterms
	\item Final: 40\%
\end{enumerate}

\section{What's Linalg?}
\vocab{Linear}: Multiplication and addition, generally.

\vocab{Algebra}: Formal structure of math, including weird things like Abstract Algebra
(and I'm pretty sure there's gonna be some group theory sprinkled in because HONORS FIFTY FOUR)

\subsection{Basic Definitions} % did i lie? did i lie? did i fucking lie?
\begin{itemize}
	\item Groups, incl. Linear Group. Can be added and subtracted from with closure, invertibility, 
	associativity, and the existence of some identity value.
	\item Rings, incl. Polynomial Ring. (set, addition, multiplication, basically the same
	properties as a group except you don't need to invert multiply entirely
	\item Fields, where you can invert multiplication entirely as well.
\end{itemize}

\vocab{Vector}:
\begin{itemize}
	\item Something w/ magnitude and direction
	\item A column of numbers
	\item An element of a vector space
	\item An ordered collection of "magnitude".
\end{itemize}

\newpage
For vectors on a plane:
\begin{itemize}
	\item Directed segment $\vec{AB}$.
	\item Translations don't change the vector. If $\vec{AB}$ is translated to $\vec{A'B'}$, then 
	$\vec{AB} = \vec{A'B'}$.
	\begin{itemize}
		\item More formally, a vector is an equivalence class of directed segments modulo translation.
		\item So $\vec{AB}$ and $\vec{A'B'}$ are different as dir. seg. and equal as vectors
	\end{itemize}
\end{itemize}

Notation: $[\vec{AB}]$ is the vector, whereas $\vec{AB}$ is just the segment.

\subsection{Vector Operations}
\begin{itemize}
	\item \textbf{addition:} $[\vec{AB}], [\vec{CD}]$ are two vectors. $[\vec{AB}]+[\vec{BE}] = [\vec{AE}]$. 
	\begin{lemma}[Commutativity of vectors]
		$[\vec{AB}] + [\vec{CD}] = [\vec{CD}] + [\vec{AB}]$.
	\end{lemma}
	\begin{proof}
		Let $E$ be defined such that $\vec{AE} \parallel \vec{CD}$ and let $F$ be defined such that 
		$\vec{BF} \parallel \vec{CD}$. Then, $\vec{EF} \parallel \vec{AB}$, meaning that the polygon
		formed by $A, B, E, F$ is a parallelogram. Letting $[\vec{AB}] = [\vec{EF}] = \vec u$ and
		$[\vec{AF}] = [\vec{BE}] = \vec v$, we see that $\vec u + \vec v = \vec v + \vec u = [\vec{AF}]$.
	\end{proof}
	\item \textbf{scalar multiplication:} $3 \cdot \vec v = \vec v + \vec v + \vec v$. We can cut a third
	of $\vec v$ out of $\vec v$, and that segment would equal $\frac{1}{3} \vec v$.
	\item \textbf{linear combinations:} If $\vec v, \vec w$ are vectors and $a, b \in \RR$, then
	$a \cdot \vec v + b \cdot \vec w$ is well-defined.
	\item \textbf{inner (dot) product:} given two vectors $\vec v, \vec u$ we define $\langle\vec v, \vec u\rangle \in \RR$.
	$\langle \vec v, \vec u \rangle = |\vec u| |\vec v| \cos \theta$. 
	\begin{itemize}
		\item The inner product is associative. 
		\item $\langle \vec v_1 + \vec v_2, \vec w \rangle = \langle \vec v_1, 
		\vec w \rangle + \langle \vec v_2, \vec w \rangle$.
		\item $\langle \vec v_1 + \vec v_2, \vec w_1 + \vec w_2 \rangle = \langle \vec v_1, 
		\vec w_1 \rangle + \langle \vec v_2, \vec w_1 \rangle + \langle \vec v_1, \vec w_2 \rangle + 
		\langle \vec v_2, \vec w_2 \rangle$.
		\item $\langle 3 \vec v, \vec w \rangle = 3 \langle \vec v, \vec w \rangle$.
		\item So the inner product is a bilinear operation as it keeps linear properties in both inputs
		and it maps everything into a third vector space.
	\end{itemize}
\end{itemize}

\subsection{Vector Spaces}
A set that supports addition into itself and scalar multiplication. We can call a vector space an 
Abelian group because vector addition is commutative and vector spaces necessarily satisfy all other 
conditions.

\subsection{Basis Vector}
Most notable example is $\vec i, \vec j, \vec k$. All other vectors in a vector space can be made with a 
linear combination of the bases.
\begin{lemma}
	$V = \langle V, e_1 \rangle e_1 + \langle V, e_2 \rangle e_2$ for a 2-dimension space. 
\end{lemma}
\begin{proof}
	Suppose $V = a_1 e_1 + a_2 e_2$. Then we apply $\langle \cdot, e_1 \rangle$ to both 
	$\langle v, e_1 \rangle = \langle a_1e_1 + a_2e_2, e_1 \rangle = a_1 \langle e_1, e_1 \rangle + a_2
	\langle e_2, e_1 \rangle$. Since $\langle e_2, e_1 \rangle$ must be $0$ since $e_1, e_2$ are both
	bases, then $V$ must equal $a_1 e_1$. The same logic can be used for the $a_2 e_2$ side. 
\end{proof}

In general, if two vectors $E_1, E_2 \in \RR^2$ satisfy the property that any vector $v \in \RR^2$ can be 
written as $V = a_1E_1 + a_2E_2$, then we say $(E_1, E_2)$ is a basis. In quantum, one notable basis is $\left \{ \begin{bmatrix}
 \frac{1}{\sqrt 2} \\ \frac{1}{\sqrt 2} \end{bmatrix}, 
\begin{bmatrix} \frac{1}{\sqrt 2} \\ -\frac{1}{\sqrt 2} \end{bmatrix} \right \}$ or $\{ \ket +, \ket - \}$.

In a basis $E_1, E_2$, we say a vector $V$ has coordinate $(a_1, a_2)$ if $V = a_1E_1 + a_2E_2$. If the basis changes, the
 coordinates also change.
\end{document}